\section{Discussion}
\label{sec:discussion}

\subsection{Why the contract is an intelligent-interface contribution}

The contract does not generate the linguistic content that makes an agent
appear intelligent.  It determines whether the interface may use and present
that content for a particular spatial interaction.  This distinction matters
because a wrong-target answer can be both fluent and locally plausible.  By
binding the pending selection to one provider and action before generation, the
system turns a silent semantic-looking failure into a deterministic admission
decision.

The user-facing consequence is limited but concrete.  A deictic request cannot
proceed from a missing or ambiguous selection; a returned answer cannot be
presented as grounded when its target or provider disagrees with the pending
interaction.  The current implementation exposes the detailed route and
binding primarily to developers, not visitors.  A future interface could use
the same evidence to explain, for example, ``This object is handled by the
Reception Agent'' or to offer a clarification choice.  Whether such
explanations improve understanding, reliance, or recovery remains an empirical
question.

\subsection{What the direct-wiring comparison establishes}

The comparison is intentionally less dramatic than an architecture
competition.  For explicit asset ownership, handoff edges, and a fixed priority
rule, direct artifacts can produce the same routing decisions and detect the
same three common fault types.  This is an important result: introducing
\systemname{} did not change the behavior already encoded in the source
artifacts, but neither did the migration automatically improve it.

The added value appears where the interaction crosses artifact boundaries.
Direct wiring has no typed occurrence connecting a scenario target to one
provider, capability, binding, action, and assertion.  Consequently, relations
such as ``this capability use has no provider'' or ``this domain capability is
provided by the runtime orchestrator'' are not representable as baseline
invariants.  The strict provider contract detected all such injected faults.
This supports a scoped adoption argument: the model is useful when a system
needs lifecycle-wide target--provider--action evidence, not merely an owner
lookup table.

\subsection{When the added structure is justified}

A static room with one chatbot and a small, stable set of objects may not
justify more than direct configuration and tests.  The contract becomes more
plausible when several agents divide responsibility, the same action kind has
many target-specific bindings, scenes are regenerated, or developers need to
trace a failure across Unity and backend artifacts.  Even then, its cost is not
free.  The current representation stores more references for the common
patches and its normalized adapter is substantially slower than direct wiring,
although the measured workload remains small relative to a typical model call.

The benchmark does not measure authoring effort.  Fewer changed top-level files
could reduce search and coordination, while more explicit references could
increase editing burden.  A developer study must test this trade rather than
inferring it from patch counts.

\subsection{Interaction design opportunities}

The present boundary supports several interaction patterns that are not yet
implemented as evaluated visitor experiences:

\begin{description}
  \item[Clarification.] When several candidates or owners exist, the interface
  can present the alternatives instead of only blocking the request.
  \item[Transparent handoff.] Before delegating, the active agent can name the
  responsible provider and reason, then let the visitor confirm.
  \item[Correction.] A visitor or author can replace the selected target or
  report that the proposed responsibility is wrong; the model relation can then
  be revised through a typed transaction.
  \item[Abstention with location.] A rejection can distinguish ``I cannot
  identify the object,'' ``no agent is responsible,'' and ``the responsible
  agent is not reachable'' rather than presenting one generic failure.
\end{description}

These patterns use evidence the system already produces, but their wording,
timing, and cognitive cost require design and evaluation.  They are stronger
future directions than adding more internal trace fields without an interface
that uses them.

\subsection{Priority evidence needed beyond this study}

Two additions would most improve external validity.  First, a fourth
independently authored, frozen scene should contain natural group- and
zone-level responsibility, at least one unassigned object, and one genuine
ambiguity.  Resolver and validator code should be frozen before the case is
revealed.  This would test transfer without deriving every interesting
condition synthetically.

Second, expected routes should be supplied by an independent oracle or a
separately implemented resolver rather than the same shared classification
function.  A declarative table reviewed before execution would be sufficient
for a modest held-out case.  These technical additions address the strongest
circularity concern without making unsupported human claims.

A developer task study is the most relevant human follow-up if ethics review,
participants, and sufficient time are available.  Participants could diagnose
and repair target, provider, and handoff faults using direct artifacts or the
contract representation.  Completion rate, time, inspected artifacts,
incorrect repairs, and explanation accuracy would test the proposed diagnostic
benefit.  A small convenience study using only subjective ratings would add
less evidence than a well-controlled held-out technical evaluation.

\subsection{Authorization and policy composition}

``Permitted'' in this paper means declared by the interaction model.  The
contract neither authenticates a visitor nor grants access to protected data or
actuators.  A deployment can first resolve the exact target--provider--action
tuple and then submit it to an external authorization policy.  Capability
lifecycle governance can separately decide which implementation version is
active.  The current prototype implements neither identity management nor
resource authorization and should not be presented as a security boundary by
itself.

\subsection{Privacy, bias, and accountability}

A spatial-agent request can reveal utterances, selected objects, locations, and
interaction histories.  The backend receives the utterance, stable target
identifiers, ray metadata, modality, and bounded history.  Persistent runtime
events omit the utterance text in the evaluated path, but the prototype lacks
visitor-facing consent, retention, export, and deletion controls.  A
deployment should minimize stored spatial data and separate technical entity
identifiers from personal identities.

Explicit agent responsibility can make organizational assumptions easier to
inspect, but it can also formalize biased or inappropriate authority.  A
validator can establish that one provider is declared consistently; it cannot
establish that the provider assignment is fair, representative, or safe.
Model review, knowledge provenance, refusal behavior, and domain safeguards
remain necessary.
